\documentclass[12pt]{article}
\date{}
\usepackage{graphicx}
%\graphicspath{ {./images/} }
\usepackage{wrapfig}

\usepackage{lipsum}

\usepackage{ucs}
\usepackage[utf8x]{inputenc}
\usepackage[greek,english]{babel}
\newcommand{\en}{\selectlanguage{english}}
\newcommand{\gr}{\selectlanguage{greek}}


\begin{document}
\title{\en Resume \gr (Βιογραφικό)}  \maketitle



%\includegraphics [width= 0.15 \textwidth]{photo.png}   

%\gr \begin{wrapfigure}{l}{0.15\textwidth}
%\centering
%\includegraphics[width=0.125\textwidth]{photo.png}
%\end{wrapfigure}

\begin{wrapfigure}{l}{0.18\textwidth}
  \centering
  \includegraphics[width=0.980\linewidth]{photo.png}
\end{wrapfigure}

\gr Ο Κιουρτ-Μεχμεταλή Αχμέτ είναι προπτυχιακός φοιτητής (έτος εισαγωγής: 2020) του τμήματος Μηχανικών Ηλεκτρονικών Υπολογιστών και Πληροφορικής του Πανεπιστημίου Πατρών. Εθελοντής σε πολλές εργασίες όπως με τίτλο «Εξέλιξη Πολιτισμού» στη δευτεροβάθμια εκπαίδευση για να αποκτήσει δεξιότητες επικοινωνίας. Γνωρίζει σε πολύ καλό επίπεδο Αγγλικά και Τουρκικά (με πτυχία). Με υψηλά κίνητρα ενδιαφέρεται για την εξέλιξη της τεχνολογίας ιδιαίτερα με το υλικό \en(hardware), \gr αλλά και με  \en database systems. \gr Έχει εμπειρία με εργασίες προγραμματισμού σε επίπεδο λογισμικού όπως σε γλώσσες \en C,C++,SQL,Java.\gr  Στο υλικό συγκεκιρμένα απασχολείτε με \en IC (Integrated Circuit), Verilog \gr και \en System Verilog.


 \bigskip

\begin{wrapfigure}{l}{0.18\textwidth}
  \centering
  \includegraphics[width=0.980\linewidth]{image2.png}
\end{wrapfigure}

\gr Ο Θεμιστοκλής Σταμαπόπουλος είναι προπτυχιακός φοιτητής του τμήματος Μηχανικών Ηλεκτρονικών Υπολογιστών και Πληροφορικής του Πανεπιστημίου Πατρών. Έχει δεξιότητες στο προγραμματισμό όπως σε \en C/C++, Java, x86 Assembly, JavaScript, HTML \gr και \en CSS. \gr Γνωρίζει την χρήση των λειτουργικών συστημάτων όπως \en Windows, Linux \gr και \en Ubuntu. \gr Έχει εργασιακή  εμπειρία ως \en Ligua Expert \gr στο τμήμα εκμάθησης Αγγλικών και Πληροφορικής για θέμα διδακτική πείρα στα μαθήματα Αγγλικών και Πληροφορικής σε μαθητές δημοτικού, γυμνάσιου και λυκείου όπως και για το θέμα ανάπτυξη εφαρμογών σε διάφορα περιβάλλοντα και λειτουργικά συστήματα \en(Windows, Ubuntu)

\end{document} 

%\begin{figure}
%\includegraphics{photo.png}
%\end{figure}


